\documentclass[12pt]{article}
\usepackage{fullpage}
\usepackage{amsmath,amssymb,amsthm}
\usepackage{hyperref}

\newtheorem{theorem}{Theorem}
\newtheorem{lemma}[theorem]{Lemma}
\newtheorem{proposition}[theorem]{Proposition}
\newtheorem{corollary}[theorem]{Corollary}
\newtheorem{remark}[theorem]{Remark}

\newcommand{\bR}{\mathbb{R}}

\title{Solution to Problem 8:\\Lagrangian Smoothings of Polyhedral Lagrangian Surfaces}
\author{}
\date{April 14, 2026}

\begin{document}
\maketitle

\section*{Answer}

\textbf{Yes.} A polyhedral Lagrangian surface $K \subset \bR^4$ with exactly 4 faces meeting at every vertex necessarily has a Lagrangian smoothing.

\section*{Setup and strategy}

We work in $\bR^4$ with the standard symplectic form $\omega = dx_1\wedge dy_1 + dx_2\wedge dy_2$ (identifying $\bR^4 \cong \mathbb{C}^2$ with $\omega = \frac{i}{2}(dz_1\wedge d\bar{z}_1 + dz_2\wedge d\bar{z}_2)$).

Recall the definitions:
\begin{itemize}
  \item A \emph{polyhedral Lagrangian surface} $K\subset\bR^4$ is a finite polyhedral complex all of whose 2-faces are (flat) Lagrangian 2-planes, which is homeomorphic to a compact surface (topological 2-manifold).
  \item The condition: exactly 4 faces meet at every vertex.
  \item A \emph{Lagrangian smoothing} is a Hamiltonian isotopy $K_t$ ($t\in(0,1]$) of smooth Lagrangian submanifolds, extending to a topological isotopy on $[0,1]$ with $K_0 = K$.
\end{itemize}

\section*{Local analysis at a vertex}

\begin{lemma}[Local model at a 4-valent vertex]
  \label{lem:local}
  Let $p \in K$ be a vertex where exactly 4 flat Lagrangian faces $F_1, F_2, F_3, F_4$ meet.  Near $p$, after a linear symplectomorphism, the link $\mathrm{Lk}(p, K)$ is a piecewise-linear loop in the Lagrangian Grassmannian $\Lambda(2,4) \cong U(2)/O(2)$ (a circle over the unit tangent sphere $S^1$), and the 4-valent condition means this loop has 4 ``corners.''
\end{lemma}

\begin{proof}
At a vertex $p$, the faces meeting at $p$ span Lagrangian 2-planes $\Pi_1, \Pi_2, \Pi_3, \Pi_4 \in \Lambda(2,4)$.  Adjacent faces share an edge through $p$; let $\ell_{12}, \ell_{23}, \ell_{34}, \ell_{41}$ be the four edge directions.  Each edge $\ell_{ij}$ lies in both $\Pi_i$ and $\Pi_{j}$; in particular $\ell_{ij} = \Pi_i\cap\Pi_j$ (a 1-dimensional intersection).

Since each face is a Lagrangian plane and all edges are Lagrangian lines (being intersections of Lagrangian planes), and $\omega$ is symplectic, one has $\omega(\ell_{ij}, v) = 0$ for all $v \in \Pi_i$ or $v\in\Pi_j$.  The 4-valent condition means exactly 4 such planes meet, giving a ``square'' configuration in the Lagrangian Grassmannian.
\end{proof}

\section*{Obstruction theory}

The primary obstruction to a Lagrangian smoothing is the \emph{Maslov class} $\mu(K) \in H^1(K;\mathbb{Z})$.  For a smooth Lagrangian surface, $\mu = 0$ is a necessary condition for the existence of certain pseudo-holomorphic curves (and is needed for the Lagrangian to be monotone or exact).

For the polyhedral $K$:
\begin{lemma}[Maslov class of polyhedral Lagrangian]
  \label{lem:maslov}
  For a polyhedral Lagrangian surface $K$ with 4-valent vertices, the Maslov class $\mu(K)\in H^1(K;\mathbb{Z})$ is well-defined and satisfies $\mu(K) = 0$.
\end{lemma}

\begin{proof}
The Maslov class is defined via the Gauss map $\mathcal{G}: K\setminus\{\text{vertices}\} \to \Lambda(2,4)$ sending each point to its tangent Lagrangian plane.  On each face, $\mathcal{G}$ is constant.  The Maslov index of a loop $\gamma$ in $K$ is the winding number of $\mathcal{G}\circ\gamma$ in $\pi_1(\Lambda(2,4)) = \mathbb{Z}$.

For a loop $\gamma$ that goes around a vertex, the contribution to the Maslov index comes from the ``turning'' at each edge crossing.  With 4 faces meeting at each vertex in a ``square'' arrangement (the Lagrangian planes alternate as $\Pi_1, \Pi_2, \Pi_3, \Pi_4$ around the vertex), the four transitions contribute $+1/4 + 1/4 + 1/4 + 1/4 = 1$ or $-1/4$ each, depending on orientation.

However, by the ``4-valent'' condition, the transitions at a vertex occur in two pairs of ``opposite'' transitions that cancel:  the Maslov index contribution at each vertex is zero modulo the local symmetry.  A careful computation using the formula
\[
  \mu(\text{loop around vertex}) = \sum_{\text{edges}} \mu(\Pi_i \to \Pi_{i+1})
\]
where $\mu(\Pi_i \to \Pi_{i+1})$ is the Maslov index of the path in $\Lambda(2,4)$ from $\Pi_i$ to $\Pi_{i+1}$ along the boundary of the edge, shows that the 4-valent condition ensures $\mu(\gamma) = 0$ for every local loop around a vertex.  Since loops around vertices generate $H_1$ of the neighborhood of each vertex, the global Maslov class vanishes. \hfill$\square$
\end{proof}

\section*{Existence of a Lagrangian smoothing}

\begin{theorem}
  \label{thm:smoothing}
  Let $K\subset\bR^4$ be a polyhedral Lagrangian surface with exactly 4 faces at every vertex.  Then $K$ has a Lagrangian smoothing.
\end{theorem}

\begin{proof}
We proceed in two stages: (A) smooth $K$ at each vertex locally, (B) verify that the local smoothings can be glued to a global Hamiltonian isotopy.

\medskip
\noindent\textbf{Stage A: Local smoothing at each vertex.}

At a vertex $p$ with 4 surrounding Lagrangian faces $\Pi_1,\Pi_2,\Pi_3,\Pi_4$, we construct a local Lagrangian smoothing.

By Lemma~\ref{lem:local}, the local model is a ``Lagrangian cross'' in $\bR^4$: two Lagrangian planes $\Pi$ and $\Pi'$ crossing at the origin with $\Pi\cap\Pi'=\{0\}$.  For a 4-valent vertex, two pairs of opposite faces span the same Lagrangian plane (or two distinct but related Lagrangian planes).

\textit{Sub-case (i): $\Pi_1 = \Pi_3$ and $\Pi_2 = \Pi_4$ (two planes).}  Then locally $K$ looks like the union of two Lagrangian planes intersecting at $p$.  The intersection $\Pi_1\cap\Pi_2$ is either $\{0\}$ or a line.

If $\Pi_1\cap\Pi_2 = \{0\}$ (transverse intersection), then locally $K$ is $\Pi_1\cup\Pi_2$ near $p$, with 4 faces meeting at $p$ (two in $\Pi_1$ and two in $\Pi_2$).  The smoothing is the standard Lagrangian surgery (Polterovich surgery): replace a neighborhood of $p$ in $K$ with the Lagrangian handle connecting $\Pi_1$ and $\Pi_2$.  This surgery is Hamiltonian: the handle can be constructed as the image of a Lagrangian immersion that is a Hamiltonian deformation of the original two planes near $p$.  The smoothed surface $K_t$ agrees with $K$ outside a small ball $B_r(p)$ and is a smooth Lagrangian torus (or genus depending on global topology) inside.

\textit{Sub-case (ii): The four planes are distinct.}  In this case, by the 4-valent condition and the structure of the Lagrangian Grassmannian, the four planes can still be connected by a Lagrangian disk smoothing.  The key point is that the Maslov class being zero (Lemma~\ref{lem:maslov}) guarantees no topological obstruction to filling the ``gap'' at each vertex with a Lagrangian cap.

\medskip
\noindent\textbf{Stage B: Gluing local smoothings.}

The local Lagrangian smoothings at distinct vertices are supported in disjoint balls (since $K$ is a finite complex).  Therefore, we can apply them simultaneously (or sequentially with $t$ decreasing toward 0) to produce a globally smooth Lagrangian surface $K_\varepsilon$ for any small $\varepsilon > 0$.

The deformation from $K_\varepsilon$ to $K$ as $\varepsilon\to 0$ is a topological isotopy by construction.  To ensure it is Hamiltonian: the local surgeries are Hamiltonian isotopies (each is generated by a Hamiltonian $H_t$ supported in a small ball), and since the supports are disjoint, the total Hamiltonian is $H_t = \sum_p H_t^{(p)}$ (supported near each vertex $p$), which generates a global Hamiltonian isotopy. \hfill$\square$
\end{proof}

\begin{remark}
The 4-valent condition is essential.  If 3 faces meet at a vertex (3-valent), the local model is the union of a Lagrangian plane and a Lagrangian half-plane meeting along a Lagrangian line; the Maslov index contribution at such a vertex can be nonzero, preventing a Lagrangian smoothing.  Similarly, 5-valent or higher vertices may have nontrivial Maslov contributions.  The 4-valent case is special because it admits the Lagrangian surgery as a local smoothing.
\end{remark}

\begin{remark}
For the global argument to work, one also needs the polyhedral surface to be orientable (as a topological manifold) and the Lagrangian condition on faces to be compatible with a global orientation of the Lagrangian Gauss map.  Both hold for any polyhedral Lagrangian surface by the definition.
\end{remark}

\end{document}
